\documentclass[../DoAn.tex]{subfiles}
\begin{document}

Chương 5 đã đánh giá các thành phần mô hình và rút ra cấu hình triển khai gồm kiến trúc CNN 1D, sử dụng nhánh quán tính ở cấu hình toàn bộ ba hoạt động. Chương 6 trình bày quá trình hiện thực hóa cấu hình đó thành một ứng dụng xác thực liên tục chạy trên Android. Phần \ref{sec:apptongquan} giới thiệu kiến trúc tổng thể của ứng dụng. Phần \ref{sec:enroll} mô tả luồng đăng ký chủ sở hữu. Phần \ref{sec:service} trình bày dịch vụ xác thực liên tục trong nền cùng máy trạng thái quyết định. Phần \ref{sec:fallback} mô tả cơ chế xác thực dự phòng bằng cử chỉ lắc. Phần \ref{sec:adaptive} trình bày cơ chế cập nhật hồ sơ thích ứng. Phần \ref{sec:webdemo} giới thiệu công cụ minh họa trên nền web trình bày năng lực đa phương thức của hệ thống.

\section{Kiến trúc ứng dụng xác thực}
\label{sec:apptongquan}

Ứng dụng \textit{BioAuth Authenticator} hiện thực hóa toàn bộ đường quyết định đã chọn ở Chương 5 trên thiết bị, không phụ thuộc vào máy chủ. Ứng dụng được tổ chức thành các tầng chức năng tách biệt. Tầng suy luận chịu trách nhiệm chạy mô hình và tính điểm tin cậy, gồm bộ thông dịch mô hình, bộ thu cửa sổ cảm biến, bộ tính điểm theo đối sánh embedding và bộ tổng hợp điểm theo thời gian. Tầng dịch vụ duy trì quá trình xác thực chạy nền liên tục. Tầng giao diện đảm nhận luồng đăng ký chủ sở hữu và đăng ký cử chỉ dự phòng. Tầng dự phòng xử lý việc phát hiện và xác minh cử chỉ lắc. Cách phân tầng này tách biệt rõ phần logic suy luận với phần vòng đời hệ thống Android, giúp hệ thống dễ kiểm thử và bảo trì.

Mô hình CNN 1D sau khi xuất sang định dạng TensorFlow Lite được đóng gói trong ứng dụng cùng các tham số phụ trợ gồm bộ chuẩn hóa đặc trưng và cấu hình suy luận. Khi vận hành, bộ thông dịch nhận cửa sổ tín hiệu quán tính đã chuẩn hóa và trả về embedding, sau đó điểm tin cậy được tính bằng độ tương đồng cosine trung bình giữa embedding với tập anchor của chủ sở hữu, rồi được chuẩn hóa theo phân bố điểm của một nhóm người dùng nền (cohort) đóng gói sẵn trong ứng dụng và đưa qua hàm sigmoid, đúng theo hàm quyết định \textit{cos\_znorm} đã chọn ở Chương 5.

Khi khởi tạo lần đầu, ứng dụng cho phép người dùng chọn chế độ mô hình tương ứng hai cấu hình ngữ cảnh đã đánh giá ở Chương 5, gồm chế độ chỉ đi bộ và chế độ toàn bộ hoạt động, như minh họa trong Hình~\ref{fig:app_mode}. Cấu hình toàn bộ hoạt động được dùng làm mặc định theo kết luận ở Chương 5; mỗi chế độ gắn với một hồ sơ riêng nên việc đổi chế độ sau khi đã đăng ký sẽ yêu cầu đăng ký lại.

\begin{figure}[H]
    \centering
    \includegraphics[width=0.28\textwidth,height=0.42\textheight,keepaspectratio]{Hinh_ve/Demo_app/screen_select_mode.jpg}
    \caption{Màn hình chọn chế độ mô hình của ứng dụng xác thực}
    \label{fig:app_mode}
\end{figure}

\section{Luồng đăng ký chủ sở hữu}
\label{sec:enroll}

Trước khi hệ thống có thể xác thực, chủ sở hữu cần thực hiện đăng ký để tạo hồ sơ định danh. Trong quá trình này, ứng dụng thu lần lượt 20 cửa sổ tín hiệu quán tính từ chủ sở hữu, mỗi cửa sổ dài bốn giây nên tổng thời gian đăng ký vào khoảng 80 giây, đưa từng cửa sổ qua bộ mã hóa để sinh embedding và lưu lại toàn bộ làm tập anchor. Số lượng 20 anchor được chọn để cân bằng giữa độ bao phủ và trải nghiệm đăng ký: một tập anchor đủ lớn giúp hồ sơ bao quát được sự biến thiên tự nhiên trong cách cầm và di chuyển thiết bị của chủ sở hữu, qua đó giảm tỉ lệ từ chối nhầm, trong khi tổng thời gian thu khoảng 80 giây vẫn đủ ngắn để người dùng không thấy phiền khi đăng ký lần đầu. Tập anchor này chính là hồ sơ của chủ sở hữu, đóng vai trò mốc tham chiếu để so sánh trong quá trình xác thực. Hình~\ref{fig:app_enroll} minh họa màn hình đăng ký ở trạng thái sẵn sàng và trong khi thu mẫu.

\begin{figure}[H]
    \centering
    \begin{subfigure}[t]{0.3\textwidth}
        \centering
        \includegraphics[width=\textwidth,height=0.4\textheight,keepaspectratio]{Hinh_ve/Demo_app/screen_train.jpg}
        \caption{Sẵn sàng đăng ký}
        \label{fig:app_enroll_a}
    \end{subfigure}
    \hfill
    \begin{subfigure}[t]{0.3\textwidth}
        \centering
        \includegraphics[width=\textwidth,height=0.4\textheight,keepaspectratio]{Hinh_ve/Demo_app/screen_training.jpg}
        \caption{Đang thu mẫu}
        \label{fig:app_enroll_b}
    \end{subfigure}
    \caption{Màn hình đăng ký sinh trắc học của chủ sở hữu}
    \label{fig:app_enroll}
\end{figure}

Điểm đáng chú ý của thiết kế là việc đăng ký chỉ đòi hỏi thu một số cửa sổ trong thời gian ngắn và không cần huấn luyện lại mô hình mạng học sâu trên thiết bị, nhờ hàm quyết định dựa trên đối sánh embedding thay vì phân loại. Đặc điểm này cho phép một mô hình dùng chung được phân phối sẵn trong ứng dụng, trong khi mỗi người dùng tự tạo hồ sơ riêng một cách nhanh chóng, phù hợp với điều kiện triển khai thực tế.

\section{Dịch vụ xác thực liên tục}
\label{sec:service}

\subsection{Cơ chế thu mẫu trong nền}

Quá trình xác thực liên tục được đảm nhận bởi một dịch vụ nền hoạt động theo kiểu duy trì để không bị hệ thống thu hồi khi ứng dụng không ở tiền cảnh. Khác với ứng dụng thu thập dữ liệu vốn ghi liên tục theo phiên, dịch vụ xác thực lấy mẫu theo sự kiện nhằm tiết kiệm năng lượng: một lượt thu được kích hoạt khi màn hình được bật hoặc khi cảm biến phát hiện chuyển động đáng kể, đồng thời có thêm các lượt thu định kỳ để bảo đảm giám sát thường xuyên. Mỗi lượt thu tạo ra một cửa sổ tín hiệu được đưa vào đường suy luận để cập nhật điểm tin cậy.

\subsection{Máy trạng thái quyết định}

Điểm tin cậy của từng cửa sổ được làm mịn theo thời gian bằng trung bình động trọng số mũ trên năm điểm gần nhất, thống nhất với cơ chế đã dùng khi đánh giá. Điểm sau khi làm mịn được ánh xạ thành một trong ba trạng thái thông qua hai ngưỡng. Khi điểm đạt từ ngưỡng tin cậy trở lên, hệ thống ở trạng thái tin cậy và không can thiệp. Khi điểm rơi vào vùng trung gian, hệ thống chuyển sang trạng thái cảnh báo và tăng cường giám sát. Khi điểm giảm xuống dưới ngưỡng thấp, hệ thống chuyển sang trạng thái nghi ngờ người lạ và kích hoạt cơ chế xác thực dự phòng. Hai ngưỡng này được đặt lấy điểm cân bằng sai số xác định trên tập đánh giá ở Chương 5 làm mốc tham chiếu, sau đó nới thành một dải để tạo vùng cảnh báo: ngưỡng tin cậy đặt quanh mốc cân bằng để phần lớn thao tác hợp lệ được chấp nhận, còn ngưỡng thấp đặt dưới mốc đó để chỉ kích hoạt dự phòng khi điểm thực sự thấp. Trạng thái hiện tại được phản ánh qua thông báo thường trực của dịch vụ nền, giúp người dùng biết được tình trạng giám sát.

Việc dùng ba trạng thái với hai ngưỡng thay vì một ngưỡng nhị phân giúp hệ thống có một vùng đệm cảnh báo, tránh kích hoạt cơ chế dự phòng quá sớm chỉ vì một vài cửa sổ nhiễu, qua đó cân bằng giữa bảo mật và trải nghiệm người dùng. Hình~\ref{fig:app_main} minh họa giao diện vận hành, hiển thị trạng thái xác thực hiện tại và điểm tin cậy ở thanh trên cùng trong khi người dùng thao tác bình thường.

\begin{figure}[H]
    \centering
    \includegraphics[width=0.28\textwidth,height=0.42\textheight,keepaspectratio]{Hinh_ve/Demo_app/screen_main.jpg}
    \caption{Giao diện vận hành hiển thị trạng thái xác thực và điểm tin cậy}
    \label{fig:app_main}
\end{figure}

\section{Cơ chế xác thực dự phòng bằng cử chỉ lắc}
\label{sec:fallback}

Khi hệ thống chuyển sang trạng thái nghi ngờ người lạ, thay vì chỉ khóa màn hình, ứng dụng kích hoạt màn hình xác thực dự phòng yêu cầu người dùng thực hiện lại cử chỉ lắc cá nhân đã đăng ký trước. Cử chỉ được thiết kế dưới dạng một mã số bí mật gồm từ bốn đến tám chữ số: mỗi loạt lắc liên tiếp được bộ phát hiện cử chỉ quy đổi thành một chữ số tương ứng với số lần lắc trong khoảng từ một đến chín, người dùng dừng khoảng một giây để chuyển sang chữ số kế tiếp, và chuỗi các chữ số tạo thành một mã tương tự mã PIN nhưng được nhập bằng chuyển động thay vì chạm phím. Quá trình nhận diện dựa trên dữ liệu gia tốc kế kết hợp thuật toán phát hiện đỉnh, không yêu cầu phần cứng bổ sung. Hình~\ref{fig:app_fallback} minh họa màn hình đăng ký mẫu lắc và màn hình xác thực dự phòng khi phát hiện người lạ.

\begin{figure}[H]
    \centering
    \begin{subfigure}[t]{0.3\textwidth}
        \centering
        \includegraphics[width=\textwidth,height=0.4\textheight,keepaspectratio]{Hinh_ve/Demo_app/screen_register_fallback.jpg}
        \caption{Đăng ký mẫu lắc}
        \label{fig:app_fallback_a}
    \end{subfigure}
    \hfill
    \begin{subfigure}[t]{0.3\textwidth}
        \centering
        \includegraphics[width=\textwidth,height=0.4\textheight,keepaspectratio]{Hinh_ve/Demo_app/screen_fallback.png}
        \caption{Xác thực khi nghi ngờ người lạ}
        \label{fig:app_fallback_b}
    \end{subfigure}
    \caption{Cơ chế xác thực dự phòng bằng cử chỉ lắc}
    \label{fig:app_fallback}
\end{figure}

Để bảo vệ bí mật của người dùng, mã cử chỉ đăng ký không được lưu ở dạng thô mà được băm kèm một giá trị muối ngẫu nhiên trước khi lưu trữ; khi xác minh, dãy lắc quan sát được băm lại và so sánh với giá trị đã lưu. Cơ chế này đóng kín vòng lặp xác thực theo trình tự xác thực thụ động, phát hiện bất thường, xác thực chủ động và khôi phục phiên, đồng thời giữ tính nhất quán khi toàn bộ quá trình đều dựa trên dữ liệu chuyển động.

\section{Cập nhật hồ sơ thích ứng}
\label{sec:adaptive}

Hành vi của một người có thể thay đổi nhẹ theo thời gian, khiến hồ sơ anchor ban đầu dần kém đại diện. Để giảm hiện tượng này, ứng dụng bổ sung một cơ chế cập nhật hồ sơ thích ứng: khi một cửa sổ được xác thực với điểm tin cậy cao, embedding tương ứng có thể được thêm vào một bộ đệm anchor thích ứng với dung lượng giới hạn, theo cơ chế thay thế các anchor cũ nhất khi đầy. Nhờ đó, tập anchor dùng để so sánh phản ánh được cả hành vi gần đây của chủ sở hữu bên cạnh hồ sơ đăng ký ban đầu.

Cần lưu ý rằng cơ chế cập nhật thích ứng này là một thành phần của bản triển khai nhằm cải thiện trải nghiệm theo thời gian, nhưng chưa được đánh giá định lượng trong khuôn khổ thực nghiệm ở Chương 5; việc đánh giá tác động của nó tới độ chính xác là một hướng mở rộng trong tương lai.

\section{Công cụ minh họa đa phương thức trên nền web}
\label{sec:webdemo}

Bên cạnh ứng dụng on-device, đề tài xây dựng một công cụ minh họa trên nền web nhằm trực quan hóa quá trình ra quyết định xác thực ở mức phiên. Khác với ứng dụng Android vận hành theo thời gian thực trên từng cửa sổ, công cụ web nhận dữ liệu một phiên đã thu, tính điểm tin cậy ở mức phiên bằng cách lấy trung bình trên các cửa sổ, và hiển thị tường minh từng bước của đường quyết định. Mục tiêu của công cụ là giúp quan sát rõ năng lực phân biệt của mô hình tách khỏi nhiễu thời gian thực, đồng thời phục vụ trình diễn trên người dùng hoàn toàn nằm ngoài tập huấn luyện.

Điểm đáng chú ý của công cụ là nó trình bày đồng thời cả ba thành phần của hệ thống đa phương thức: điểm của nhánh quán tính được tính theo hàm quyết định đã chọn, điểm của nhánh tương tác màn hình, và điểm dung hợp của hai nhánh. Ứng với đó, công cụ hiển thị song song hai kết luận xác thực: một theo nhánh quán tính, vốn là đường quyết định của bản triển khai, và một theo điểm dung hợp. Việc đặt cạnh nhau hai cột kết luận cho phép người xem trực tiếp đối chiếu đóng góp của nhánh tương tác so với nhánh quán tính trên từng phiên cụ thể.

Cách trình bày này phục vụ đúng mục tiêu khảo sát đa phương thức của đề tài. Một mặt, nó cho thấy hệ thống được thiết kế để hỗ trợ nhiều phương thức và có khả năng dung hợp. Mặt khác, nó minh họa một cách trực quan kết luận đã rút ra ở Chương 5: trên các phiên thử nghiệm, điểm dung hợp không vượt trội so với điểm của riêng nhánh quán tính, và điểm của nhánh tương tác đối với người lạ vẫn ở mức thấp nên không làm tăng nguy cơ chấp nhận nhầm. Như vậy công cụ web không mâu thuẫn với lựa chọn triển khai đơn phương thức mà còn củng cố cơ sở cho lựa chọn đó: người xem có thể tự kiểm chứng rằng việc bổ sung nhánh tương tác qua dung hợp tuyến tính không cải thiện chất lượng quyết định trong điều kiện dữ liệu hiện có.

Ngưỡng quyết định trong công cụ được đặt cố định, calib một lần tại điểm cân bằng lỗi trên tập đánh giá và áp dụng chung cho mọi người dùng, nhất quán với nguyên tắc đặt ngưỡng đã trình bày ở Chương 5. Nhờ đó, khi trình diễn với người dùng mới chưa từng xuất hiện trong huấn luyện, công cụ phản ánh đúng kịch bản tập mở mà đề tài hướng tới.

Chương này đã trình bày quá trình hiện thực hóa cấu hình đã chọn thành một ứng dụng xác thực liên tục trên Android, bao gồm kiến trúc phân tầng, luồng đăng ký dựa trên anchor, dịch vụ nền với máy trạng thái ba mức, cơ chế xác thực dự phòng bằng dãy cử chỉ lắc và cơ chế cập nhật hồ sơ thích ứng, cùng một công cụ minh họa trên nền web trình bày đầy đủ năng lực đa phương thức của hệ thống. Toàn bộ đường suy luận chạy trực tiếp trên thiết bị, nhất quán với hàm quyết định đã thiết kế và đánh giá ở các chương trước. Chương 7 tổng kết các kết quả đạt được của đề tài, nêu những hạn chế còn tồn tại và đề xuất các hướng phát triển tiếp theo.

\end{document}